\section{Discussion}
\label{sec:discussion}

\paragraph{What the normal-mode analogy gets right.} Taken seriously as
an overdamped, thermally driven system, the coupled-oscillator picture
predicts that the \emph{lowest}-variance eigenportfolios relax fastest.
Out of sample, that ordering holds. Reversion profits and negative
autocorrelation are concentrated at the bottom of the spectrum, and the
market and sector modes at the top trend rather than revert. The
analogy is therefore a useful guide to where to look, and it inverts
the usual PCA habit of treating the leading components as the
interesting ones.

\paragraph{What it gets wrong.} The quantitative prediction, a half-life
proportional to $\lambda_k$, is not supported. The reversion we
measure is mostly a one-day effect that does not survive an extra day
of execution delay. That is more naturally read as liquidity provision
and microstructure noise \citep{lehmann1990,nagel2012} spread over many
small-variance combinations of stocks than as slow relaxation toward an
equilibrium.

\paragraph{The Marchenko--Pastur filter is the wrong filter for this
purpose.} RMT tells us which eigenvalues are distinguishable from noise,
and that is the right question for covariance cleaning
\citep{laloux1999,bunbouchaud2017}. For reversion trading, the relevant
object is the low-variance \emph{subspace}. Individual bulk eigenvectors
are unstable, but projecting onto many of them and trading them together
diversifies away their individual instability. This explains why
equal-weighting the whole bulk beats every selection rule we tried.

\paragraph{Machine learning adds little; timing adds more.} With AUCs of
0.51--0.53, the classifiers do find structure. They are significantly
better than random selection, and stability, concentration and
autocorrelation carry information. But their edge is smaller than the
diversification cost of trading only a third of the pool. Regime
information is more valuable as a \emph{timing} signal for the whole
book than as a feature for ranking modes. The regime asymmetry is
large: reversion is unprofitable in calm markets and profitable in
stressed ones. Our overlay exploits it walk-forward, but its form was
chosen after seeing this asymmetry.

\paragraph{Multiple testing.} We report every strategy and variant we
ran. Across them, the best net Sharpe ratio is below 0.2 over twelve
years, so none would survive a deflated-Sharpe adjustment for the
number of trials \citep{baileylopezdeprado2014}. The robust findings
are the cross-spectrum pattern (\cref{tab:spectrum}), which involves
no selection, and the regime asymmetry, which holds in every
specification.

\section{Limitations}
\label{sec:limitations}

\begin{itemize}
\item \textbf{Survivorship.} The universe is the current S\&P~500
  restricted to stocks with complete 2010--2024 histories (\DataAssets{}
  names). This removes delisted and acquired firms. Survivorship
  typically flatters long-only returns more than market-neutral
  reversion, but it can also understate the tail risk of a reversion
  book, which loses when a stock falls and never comes back.
  Point-in-time constituents (e.g.\ CRSP) would fix this.
\item \textbf{Execution.} Signals use closing prices and are traded at
  the same close. The execution-delay test shows that this assumption
  matters a great deal. The cost model is a flat charge per unit
  turnover, with no market impact, borrow costs or financing.
\item \textbf{Estimation.} We use sample correlations with a fixed
  252-day window. Shrinkage or rotationally invariant estimators
  \citep{bunbouchaud2017}, or eigenvector tracking across windows
  rather than re-estimating from scratch, could make individual modes
  more stable and selection more valuable.
\item \textbf{Label.} The binary profitability label discards magnitude,
  and the payoff is negatively skewed. Regression or quantile targets,
  or a direct portfolio objective, might use the weak signal more
  efficiently.
\end{itemize}

\section{Conclusion}

We set out to test whether the low-eigenvalue normal modes of the
equity correlation matrix mean-revert predictably, and whether
supervised learning and regime detection can turn that into a
statistical-arbitrage strategy. On 417 S\&P~500 stocks over 2013--2024,
out of sample, the answers are as follows. Yes, mode portfolios
mean-revert, but at the bottom of the spectrum rather than in the
Marchenko--Pastur signal modes, and mostly at a one-day horizon. Only
weakly can it be predicted which modes will revert: AUC 0.51--0.53,
better than chance and better than random selection, but not better
than holding everything. Market regimes predict \emph{when} reversion
pays far better than mode features predict \emph{which} modes pay.
After realistic costs, the best walk-forward implementation earns a
small positive Sharpe ratio, \SRMLOverB{}, that depends on a
data-informed overlay and on same-close execution. The
coupled-oscillator analogy is a good map of where reversion lives in
the spectrum. It is not, on this evidence, a source of tradable alpha
for a daily close-to-close strategy.
